\chapter{Дополнительные материалы}
\label{app:supplementary}

\section{Гиперпараметры обучения YOLOv8}

\begin{table}[H]
\centering
\caption{Полный список гиперпараметров обучения}
\label{tab:hp_full}
\begin{tabular}{lll}
\toprule
\textbf{Параметр} & \textbf{Значение} & \textbf{Описание} \\
\midrule
\multicolumn{3}{l}{\textit{Основные параметры}} \\
\texttt{model}            & yolov8n.pt & Базовая модель \\
\texttt{epochs}           & 100        & Число эпох \\
\texttt{batch}            & 32         & Размер батча \\
\texttt{imgsz}            & 640        & Размер изображения \\
\texttt{device}           & GPU        & Устройство обучения \\
\texttt{workers}          & 8          & Число воркеров загрузки \\
\texttt{pretrained}       & true       & Предобученные веса COCO \\
\midrule
\multicolumn{3}{l}{\textit{Оптимизатор}} \\
\texttt{optimizer}        & auto       & Выбор оптимизатора \\
\texttt{lr0}              & 0.01       & Начальная скорость обучения \\
\texttt{lrf}              & 0.01       & Финальный LR (доля от lr0) \\
\texttt{momentum}         & 0.937      & SGD momentum \\
\texttt{weight\_decay}    & 0.0005     & L2-регуляризация \\
\texttt{warmup\_epochs}   & 3.0        & Эпохи разогрева \\
\texttt{warmup\_momentum} & 0.8        & Momentum при разогреве \\
\texttt{warmup\_bias\_lr} & 0.1        & Bias LR при разогреве \\
\midrule
\multicolumn{3}{l}{\textit{Веса потерь}} \\
\texttt{box}              & 7.5        & Вес box loss \\
\texttt{cls}              & 0.5        & Вес classification loss \\
\texttt{dfl}              & 1.5        & Вес DFL loss \\
\midrule
\multicolumn{3}{l}{\textit{Аугментация}} \\
\texttt{hsv\_h}           & 0.015      & Изменение hue \\
\texttt{hsv\_s}           & 0.7        & Изменение saturation \\
\texttt{hsv\_v}           & 0.4        & Изменение value \\
\texttt{translate}        & 0.1        & Случайный сдвиг \\
\texttt{scale}            & 0.5        & Случайный масштаб \\
\texttt{fliplr}           & 0.5        & Горизонтальное отражение \\
\texttt{mosaic}           & 1.0        & Mosaic-аугментация \\
\texttt{erasing}          & 0.4        & Random erasing \\
\texttt{auto\_augment}    & randaugment & Стратегия аугментации \\
\texttt{close\_mosaic}    & 10         & Отключить mosaic последние N эпох \\
\midrule
\multicolumn{3}{l}{\textit{Детекция}} \\
\texttt{iou}              & 0.7        & IoU порог для NMS \\
\texttt{max\_det}         & 300        & Макс. детекций на изображение \\
\bottomrule
\end{tabular}
\end{table}

\section{Библиотеки Python}

\begin{table}[H]
\centering
\caption{Версии библиотек Python, использованных в исследовании}
\label{tab:python_env}
\begin{tabular}{lll}
\toprule
\textbf{Библиотека} & \textbf{Версия} & \textbf{Назначение} \\
\midrule
Python        & 3.13    & Базовый язык программирования \\
ultralytics   & 8.4.23+ & YOLOv8 реализация \\
torch         & 2.0+    & PyTorch (бэкенд Ultralytics) \\
kagglehub     & 1.0.0+  & Загрузка датасетов из Kaggle \\
litestar      & 2.21.1+ & REST API фреймворк \\
pillow        & 12.1.1+ & Обработка изображений \\
uvicorn       & 0.42.0+ & ASGI сервер \\
python-multipart & 0.0.22+ & Обработка multipart форм \\
\bottomrule
\end{tabular}
\end{table}

\section{Код подготовки датасета}

\begin{lstlisting}[language=Python, caption=Скрипт подготовки датасета (prepare\_dataset.py), label=lst:prepare]
import shutil
import random
from pathlib import Path
import kagglehub

CLASSES = ["first_degree", "second_degree", "third_degree"]
VAL_SPLIT = 0.2

def download_dataset() -> Path:
    path = kagglehub.dataset_download(
        "shubhambaid/skin-burn-dataset"
    )
    return Path(path)

def prepare_dataset(source_dir: Path, val_split: float = VAL_SPLIT) -> None:
    images_train = DATASET_DIR / "images" / "train"
    images_val = DATASET_DIR / "images" / "val"
    labels_train = DATASET_DIR / "labels" / "train"
    labels_val = DATASET_DIR / "labels" / "val"

    for d in [images_train, images_val, labels_train, labels_val]:
        d.mkdir(parents=True, exist_ok=True)

    image_files = list(source_dir.glob("*.jpg"))
    random.seed(42)
    random.shuffle(image_files)

    split_idx = int(len(image_files) * (1 - val_split))
    train_images = image_files[:split_idx]
    val_images = image_files[split_idx:]

    for img_path in train_images:
        label_path = img_path.with_suffix(".txt")
        if label_path.exists():
            shutil.copy(img_path, images_train / img_path.name)
            shutil.copy(label_path, labels_train / label_path.name)

    for img_path in val_images:
        label_path = img_path.with_suffix(".txt")
        if label_path.exists():
            shutil.copy(img_path, images_val / img_path.name)
            shutil.copy(label_path, labels_val / label_path.name)

    # Create data.yaml configuration
    data_yaml = DATASET_DIR / "data.yaml"
    data_yaml.write_text(f"""path: {DATASET_DIR}
train: images/train
val: images/val

names:
  0: {CLASSES[0]}
  1: {CLASSES[1]}
  2: {CLASSES[2]}
""")

    print(f"Dataset prepared: {len(train_images)} train, {len(val_images)} val")
\end{lstlisting}

\section{Код обучения модели}

\begin{lstlisting}[language=Python, caption=Скрипт обучения модели (train.py), label=lst:train_full]
from pathlib import Path
from ultralytics import YOLO

def train(
    epochs: int = 100,
    imgsz: int = 640,
    batch: int = 16,
    model_size: str = "n",
) -> Path:
    MODELS_DIR.mkdir(parents=True, exist_ok=True)
    data_yaml = DATASET_DIR / "data.yaml"

    if not data_yaml.exists():
        raise FileNotFoundError(
            f"Dataset not found at {data_yaml}. "
            "Run prepare_dataset.py first."
        )

    model = YOLO(f"yolov8{model_size}.pt")

    results = model.train(
        data=str(data_yaml),
        epochs=epochs,
        imgsz=imgsz,
        batch=batch,
        project=str(MODELS_DIR),
        name="burn_detector",
        exist_ok=True,
    )

    best_model = MODELS_DIR / "burn_detector" / "weights" / "best.pt"
    print(f"Training complete. Best model: {best_model}")
    return best_model
\end{lstlisting}

\section{REST API для инференса}

\begin{lstlisting}[language=Python, caption=REST API сервер (api.py), label=lst:api]
import io
from dataclasses import dataclass

from litestar import Litestar, get, post
from litestar.datastructures import UploadFile
from litestar.enums import RequestEncodingType
from litestar.openapi import OpenAPIConfig
from litestar.params import Body
from litestar.response import Response
from PIL import Image
from ultralytics import YOLO


@dataclass
class Detection:
    class_id: int
    class_name: str
    confidence: float
    bbox: list[float]


@dataclass
class PredictionResponse:
    detections: list[Detection]


@get("/health")
async def health() -> dict[str, str]:
    return {"status": "ok"}


@post("/predict")
async def predict(
    data: UploadFile = Body(media_type=RequestEncodingType.MULTI_PART),
) -> PredictionResponse:
    content = await data.read()
    image = Image.open(io.BytesIO(content))

    results = get_model().predict(image, verbose=False)
    detections = []

    for result in results:
        for box in result.boxes:
            detections.append(
                Detection(
                    class_id=int(box.cls[0]),
                    class_name=result.names[int(box.cls[0])],
                    confidence=float(box.conf[0]),
                    bbox=box.xyxy[0].tolist(),
                )
            )

    return PredictionResponse(detections=detections)


@post("/predict/visualize")
async def predict_visualize(
    data: UploadFile = Body(media_type=RequestEncodingType.MULTI_PART),
) -> Response[bytes]:
    content = await data.read()
    image = Image.open(io.BytesIO(content))

    results = get_model().predict(image, verbose=False)
    annotated = results[0].plot()

    img_pil = Image.fromarray(annotated[..., ::-1])
    buffer = io.BytesIO()
    img_pil.save(buffer, format="JPEG")
    buffer.seek(0)

    return Response(
        content=buffer.read(),
        media_type="image/jpeg",
    )


app = Litestar(
    route_handlers=[health, predict, predict_visualize],
    openapi_config=OpenAPIConfig(
        title="Burn Classifier API",
        version="1.0.0",
        description="Skin burn detection using YOLOv8",
    ),
)
\end{lstlisting}

\section{Пример использования API}

\begin{lstlisting}[language=bash, caption=Server launch and API request, label=lst:api_usage]
# Start server
uvicorn burn_classifier.api:app --host 0.0.0.0 --port 8000

# Send image for analysis (JSON response)
curl -X POST "http://localhost:8000/predict" \
     -F "data=@burn_image.jpg"

# Example response:
# {
#   "detections": [
#     {
#       "class_id": 1,
#       "class_name": "second_degree",
#       "confidence": 0.89,
#       "bbox": [120.5, 80.2, 340.8, 280.5]
#     }
#   ]
# }

# Get image with visualization
curl -X POST "http://localhost:8000/predict/visualize" \
     -F "data=@burn_image.jpg" \
     --output result.jpg
\end{lstlisting}

\section{Скрипт настройки vast.ai}

\begin{lstlisting}[language=bash, caption=Cloud environment setup script (vastai\_setup.sh), label=lst:vastai]
#!/bin/bash
# vast.ai GPU server environment setup

# Install uv (fast Python package manager)
pip install uv

# Create virtual environment
uv venv && source .venv/bin/activate

# Install dependencies
uv sync

# Prepare dataset
uv run prepare-dataset

# Start training
uv run train --epochs 100 --batch 32 --model-size n

# After training - start API server
uv run serve
\end{lstlisting}

\section{Структура проекта}

\begin{verbatim}
burn-classifier/
|-- src/
|   |-- burn_classifier/
|       |-- __init__.py          # Пути и константы
|       |-- prepare_dataset.py   # Подготовка данных
|       |-- train.py             # Обучение модели
|       |-- api.py               # REST API
|-- dataset/
|   |-- images/
|   |   |-- train/               # 980 изображений
|   |   |-- val/                 # 245 изображений
|   |-- labels/
|   |   |-- train/               # Аннотации YOLO
|   |   |-- val/
|   |-- data.yaml                # Конфигурация
|-- models/
|   |-- burn_detector/
|       |-- weights/
|           |-- best.pt          # Лучшая модель
|           |-- last.pt          # Последняя эпоха
|-- pyproject.toml               # Зависимости
|-- vastai_setup.sh              # Настройка облака
|-- yolov8n.pt                   # Базовые веса
\end{verbatim}
